\scanpage{429}
\mspage{341}

\uline{Démonstration} On applique 4.11.1 \inserted{(p. 267-268)} \unsure{séparément|séparement} à $\unsure{M|\mathcal M}\approx\mathrm{Ind}_\pi M\deleted{(\pi)}_\pi$ et à $S\approx\mathrm{Ind}_\pi\,\underline S(\pi)\unsure{\}|)}$, pour conclure que $M$ et $\underline S$ sont stables séparément par $\varinjlim$ de type $\unsure{\Phi|\mathscr I|\mathscr D}$. Il reste à voir que le foncteur $\underline S\longrightarrow M$ y commute. Cela se voit comme dans la démonstration de 4.11.1, \unsure{on} obtient plus généralement ceci :

\uline{Lemme} 4.13.20 (cor. à 4.11.1) : \unsure{soient|soient\textsuperscript{t}} $f_0:C\longrightarrow C'$, un foncteur entre catégories \unsure{équivalentes à|équiv. à} des petites catégories. $\pi$ un cardinal infini, $f:M\overset{\text{déf}}{=}\mathrm{Ind}_\pi C\longrightarrow M'\overset{\text{déf}}{=}\mathrm{Ind}_\pi C'$
% slanted
\margin{Il vaut mieux \unsure{prendre|parler de} un foncteur $f_0:C\longrightarrow N$, et dire que $f_0$ commutant aux $\varinjlim$ de type $\Phi$, \unsure{alors} il en est de $=$ de son \unsure{extension|extens.} $\mathrm{Ind}_\pi(C,N)\longrightarrow N$. (NB on suppose $N$ satisfait $\mathrm L_\pi$).}
l'extension canonique $\mathrm{Ind}_\pi(f_0)$ de $f_0$. Soit $\Phi$ \inserted{et $\alpha$} comme dans 4.12.1, et supposons que $\deleted{\unsure{I_\alpha}}I_\alpha\in\Phi$ et que $f_0$ commute aux $\varinjlim$ de type $\Phi$. Alors $f$ commute à ces $\unsure{\Rightarrow|=}$ limites inductives.

\uline{Corollaire} 4.13.21. \inserted{Dans les conditions de 4.13.20,} Supposons que $\Phi$ contienne les catégories de cardinal $\leq\pi$ (resp. les catégories filtrantes de cardinal $\leq\pi$, \inserted{resp. des cat. filtr. \unsure{grdes|grds} devant $\pi_0$ de card. $\leq\pi$,} resp. les catégories discrètes de cardinal $\leq\pi$). Alors $f$ commute aux \inserted{petites} $\varinjlim$ \deleted{quelconques} (resp. aux petites $\varinjlim$ filtrantes, \inserted{resp. aux petites $\varinjlim$ filtrantes grandes devant $\pi_0$,} resp. aux petites sommes).

(Réciproque claire et ``laissée au lecteur''.)

Se voit comme 4.4.14 et 4.4.14.1 (p. 97,98).

% archive sheet number: 428

